% Part: computability
% Chapter: recursive-functions
% Section: examples

\documentclass[../../../include/open-logic-section]{subfiles}

\begin{document}

\olfileid[fr]{cmp}{rec}{exa}
\olsection{Exemples de fonctions récursives primitives}

Nous connaissons déjà quelques fonctions récursives primitives :
l'addition~$\Add$ et la multiplication~$\Mult$. La fonction
identité $\fn{id}(x)=x$ est récursive primitive, puisqu'elle
n'est autre que~$\Proj{1}{0}$. Les fonctions constantes
$\fn{const}_n(x)=n$ le sont également : elles s'obtiennent
par compositions successives de $\Zero$ et de~$\Succ$.
Cela permet d'employer des constantes dans les définitions
récursives primitives. Ainsi, pour définir $f(x)=2\cdot x$,
on peut composer la fonction constante de valeur~$2$
avec la multiplication :
$f(x)=\Mult(\fn{const}_2(x),\Proj{1}{0}(x))$.
Nous utiliserons désormais cette méthode.

\begin{prop}
  La fonction d'exponentiation
  $\fn{exp}(x,y)=x^y$ est récursive primitive.
\end{prop}

\begin{proof}
  On peut définir $\fn{exp}$ par récurrence primitive :
  \begin{align*}
    \fn{exp}(x, 0) & = 1\\
    \fn{exp}(x, y+1) & = \Mult(x, \fn{exp}(x,y)).
    \intertext{À strictement parler, cette écriture ne donne pas
      encore une définition à partir de fonctions dont on sait
      qu'elles sont récursives primitives. Sous la forme requise,
      elle devient :}
    \fn{exp}(x, 0) & = f(x)\\
    \fn{exp}(x, y+1) & = g(x, y, \fn{exp}(x,y)).
    \intertext{où}
    f(x) & = \Succ(\Zero(x)) = 1\\
    g(x, y, z) & = \Mult(\Proj{3}{0}(x, y, z), \Proj{3}{2}(x, y, z)) = x \cdot z
  \end{align*}
  Les fonctions $f$ et~$g$ s'obtiennent donc par composition
  à partir de fonctions récursives primitives.
\end{proof}

\begin{prop}
  La fonction prédécesseur $\fn{pred}(y)$, définie par
  \[
  \fn{pred}(y) = \begin{cases}
    0 & \text{si $y=0$}\\
    y-1 & \text{sinon}
  \end{cases}
  \]
  est récursive primitive.
\end{prop}

\begin{proof}
 Observons que
 \begin{align*}
   \fn{pred}(0) & = 0 \text{ et}\\
   \fn{pred}(y+1) & = y.
 \end{align*}
 Ces équations ressemblent à une définition par récurrence
 primitive, mais le schéma formel exige au moins un argument
 supplémentaire~$x$. Il est également inhabituel que la
 définition de $\fn{pred}(y+1)$ n'utilise pas la valeur
 précédente~$\fn{pred}(y)$. Définissons d'abord
 $\fn{pred}'(x,y)$ par
 \begin{align*}
   \fn{pred}'(x, 0) & = \Zero(x) = 0,\\
   \fn{pred}'(x, y+1) & = \Proj{3}{1}(x, y, \fn{pred'}(x, y)) = y.
 \end{align*}
 La composition donne alors, par exemple,
 $\fn{pred}(x)=\fn{pred}'(\Zero(x),\Proj{1}{0}(x))$.
\end{proof}

\begin{prop}
  La fonction factorielle
  $\fn{fac}(x)=\fact{x}=1\cdot2\cdot3\cdot\dots\cdot x$
  est récursive primitive.
\end{prop}

\begin{proof}
  Les équations récursives naturelles sont
  \begin{align*}
    \fn{fac}(0) &= 1\\
    \fn{fac}(y+1) & = \fn{fac}(y) \cdot (y+1).
    \intertext{Pour obtenir la forme exigée, définissons d'abord
      une fonction~$h$ à deux arguments :}
    h(x, 0) & = \fn{const}_1(x)\\
    h(x, y+1) & = g(x, y, h(x, y))
    \intertext{où
      $g(x,y,z)=\Mult(\Proj{3}{2}(x,y,z),
      \Succ(\Proj{3}{1}(x,y,z)))$, puis posons}
    \fn{fac}(y) & = h(\Proj{1}{0}(y), \Proj{1}{0}(y)) = h(y,y).
  \end{align*}
  Dans la suite, nous omettrons parfois ces développements
  formels en compositions et en récurrences primitives.
\end{proof}

\begin{prop}
  La soustraction tronquée $x\tsub y$, définie par
  \[
  x \tsub y = \begin{cases}
    0 & \text{si $x<y$}\\
    x-y & \text{sinon}
  \end{cases}
  \]
  est récursive primitive.
\end{prop}

\begin{proof}
  On a :
  \begin{align*}
    x \tsub 0 & = x\\
    x \tsub (y+1) & = \fn{pred}(x \tsub y)
  \end{align*}
\end{proof}

\begin{prop}
 La distance entre $x$ et $y$, $\left|x-y\right|$,
 est récursive primitive.
\end{prop}

\begin{proof}
  Puisque $\left|x-y\right|=(x\tsub y)+(y\tsub x)$,
  la distance s'obtient par composition à partir de
  l'addition et de $\tsub$, qui sont récursives primitives.
\end{proof}

\begin{prop}
  Le maximum de $x$ et de $y$, $\fn{max}(x,y)$,
  est récursif primitif.
\end{prop}

\begin{proof}
  On peut définir $\fn{max}(x,y)$ par composition
  à partir de l'addition et de $\tsub$ :
  \[
  \fn{max}(x,y) \defis x + (y \tsub x).
  \]
  Si $x$ est le maximum, donc si $x\ge y$, on a
  $y\tsub x=0$, d'où $x+(y\tsub x)=x+0=x$.
  Si $y$ est le maximum, $y\tsub x=y-x$,
  d'où $x+(y\tsub x)=x+(y-x)=y$.
\end{proof}

\begin{prop}
  \ollabel{prop:min-pr}
  Le minimum de $x$ et de $y$, $\fn{min}(x,y)$,
  est récursif primitif.
\end{prop}

\begin{proof}
  Exercice.
\end{proof}

\begin{prob}
  Démontrer la \olref[cmp][rec][exa]{prop:min-pr}.
\end{prob}

\begin{prob}
Montrer que \[
f(x, y) =
2^{(2^{\iddots^{2^{x}}})}\raisebox{1ex}{\bigg\rbrace}
\raisebox{1ex}{\text {$y$ occurrences de $2$}}\]
est récursive primitive.
\end{prob}

\begin{prob}
Montrer que la division entière
$d(x,y)=\lfloor x/y\rfloor$, c'est-à-dire la division
où l'on néglige la partie après la virgule, est récursive
primitive. Pour $y=0$, on convient que $d(x,y)=0$.
Donner une définition explicite de~$d$ au moyen de
la récurrence primitive et de la composition.
\end{prob}

\begin{prop}
L'ensemble des fonctions récursives primitives est
clos par les deux opérations suivantes :
\begin{enumerate}
\item Sommes finies : si $f(\vec x,z)$ est récursive primitive,
la fonction suivante l'est aussi :
\[
g(\vec x, y) \defis \sum_{z = 0}^y f(\vec x, z).
\]
\item Produits finis : si $f(\vec x,z)$ est récursive primitive,
la fonction suivante l'est aussi :
\[
h(\vec x, y) \defis \prod_{z = 0}^y f(\vec x, z).
\]
\end{enumerate}
\end{prop}

\begin{proof}
Par exemple, les sommes finies sont définies récursivement par
\begin{align*}
  g(\vec x, 0) & = f(\vec x, 0)\\
  g(\vec x, y+1) & = g(\vec x, y) + f(\vec x, y+1).
\end{align*}
\end{proof}

\end{document}
